% !TeX root = ../technical_paper.tex
\section{作品难点与创新}

本作品面向户用储能、离网备用供电与远程状态监测等应用场景，围绕 GD32G553 主控、GD32C113 BMS 和 GD32VW553 无线模块构建多 MCU 协同系统。GD32G553 负责系统级调度、BMS 数据接收与转发、CC/CV 充电闭环控制及 PWM 输出；GD32C113 负责电芯采样、均衡、SOC 估算和故障保护；GD32VW553 负责无线通信和远程监控。主控侧采用 APP / Middleware / Protocol / HAL 四层解耦结构，BMS 侧以产品状态机实现测量、保护、FET、均衡和 SOC 的协同管理。论文后续章节将围绕系统架构、硬件设计、软件流程和测试验证展开。

\subsection{应用场景与作品定位}

随着户用储能和分布式能源需求的增长，电池管理系统不仅需要完成基本的数据采集和保护动作，还需要满足状态展示、远程通信和调试验证等多方面要求。本作品的定位是围绕"采样-判定-保护-通信-展示"完整链路，构建一套多 MCU 协同的数字电源管理控制平台。

与传统单控制器集中式方案相比，本作品通过主控、BMS 和无线模块的分层设计，将实时性要求高的控制和保护逻辑保留在本地，将刷新频率较低的监控任务移出核心控制路径，避免通信交互对控制周期造成干扰。

\subsection{难点分析}

实现光伏储能一体化的数字电源系统，主要面临以下难题：

\begin{enumerate}[label=\arabic*、, leftmargin=*, itemsep=0.8ex, parsep=0pt, topsep=0.5ex]
  \item \textbf{实时控制与非实时通信的解耦。}
  主控侧既要处理采样、PWM 输出和保护动作等高实时任务，又要承担 BMS 数据转发和无线通信等低实时任务。若任务在同一层级混合，易造成中断响应受阻和控制周期抖动。

  \item \textbf{主控、BMS 与无线模块的职责边界。}
  主控负责系统状态组织和数据汇总，BMS 负责电池采样、保护、均衡和 FET 策略，无线模块负责展示和交互。职责边界不清晰会导致重复判断和状态冲突，必须通过统一协议和状态码约束。

  \item \textbf{数据协议与状态定义的统一。}
  主控与 BMS 之间使用 CAN FD 遥测，主控与无线模块之间使用 UART 帧转发，若电压、电流、SOC 和告警码采用不同编码，会造成显示、日志和控制逻辑不一致。

  \item \textbf{测量与标定的准确性。}
  主控 ADC 采样和 BMS 电芯采样均需完成零点校准、增益换算和物理量转换。若标定方法不一致或采样链路存在漂移，将导致保护阈值误判和 SOC 偏差。

  \item \textbf{磷酸铁锂电池 SOC 估算精度。}
  LiFePO$_4$ 电池在 $30\%$--$80\%$ SOC 区间的 OCV 曲线较为平坦，单纯依靠开路电压法难以获得准确初值；安时积分法又存在累计误差和初值依赖问题。

  \item \textbf{三端口功率变换的协同设计。}
  三端口变换需要同时考虑光伏端、电池端和输出端之间的能量流动，控制对象比单输入单输出变换器更复杂。功率板中存在高 $dv/dt$ 开关节点，对隔离驱动、采样抗干扰和 PCB 布局提出了较高要求。
\end{enumerate}

\subsection{创新点分析}

针对上述难点，本作品从系统架构、电池管理、状态估算、功率平台和远程监控五个层面进行设计，形成如下创新。

\begin{enumerate}[label=\arabic*、, leftmargin=*, itemsep=0.8ex, parsep=0pt, topsep=0.5ex]
  \item \textbf{多 MCU 分层协同控制架构。}
  采用 GD32G553、GD32C113 和 GD32VW553 三颗 MCU 协同工作，将系统划分为主控层、电池管理层和无线交互层。主控侧采用 APP / Middleware / Protocol / HAL 四层解耦结构，通过统一调度器和发布/订阅总线组织任务；BMS 侧负责电池测量、保护、均衡和 SOC 管理；无线侧负责远程监控与参数交互。该架构将高实时性任务与非实时任务分离，提高系统稳定性。

  \item \textbf{基于产品状态机的智能 BMS 管理体系与四级安全防线。}
  GD32C113 BMS 工程采用具有明确状态边界的产品状态机，包含 INIT、RUN 和 DEGRADED 三种状态。在安全保护方面，构建了 AFE 硬件保护、MCU 软件备份保护、看门狗监控和 FETOFF 硬关断的四级安全防线，涵盖过压、欠压、过流、短路、过温等多重故障保护。

  \item \textbf{OCV-安时积分融合 SOC 估算策略。}
  针对 LiFePO$_4$ 电池平坦的 OCV-SOC 特性曲线，采用 OCV 种子自适应初始化结合动态库仑计数的混合估算策略。系统在上电或长时间静置后通过 OCV-SOC 查表获得 SOC 初值；运行过程中以安时积分法进行动态累计，并定期通过 OCV 修正消除累计漂移。

  \item \textbf{面向光储应用的三端口功率变换平台设计。}
  设计了三端口功率变换平台，完成拓扑设计、功率板设计、PCB 实现以及驱动与采样电路设计。目前该平台已完成 PV 端对储能端的充电功能验证，为后续实现 PV$\rightarrow$Load、Battery$\rightarrow$Load 以及多模式能量流控制提供了统一硬件基础。

  \item \textbf{基于 MQTT 的轻量级云端监控平台。}
  无线模块通过 MQTT 协议实现与云端的实时双向通信，支持遥测数据上报和远程命令下发。同时开发了基于 Python 的 Web 监控服务器，采用 SSE 技术实现浏览器端遥测数据实时可视化。控制链路与监控链路分离的设计保证了即使监控端出现延迟或掉线，也不会影响底层保护和运行控制。
\end{enumerate}

\subsection{本章小结}

本章围绕作品定位、实现难点和创新进行了说明。本作品通过多 MCU 分层协同、产品状态机管理、OCV-安时积分混合 SOC 估算、四级安全防线和 MQTT 云端监控，构建了数字电源管理控制平台。三端口功率变换平台已完成硬件设计与基础功能验证，为后续实现完整光储能量管理奠定了硬件基础。
